\section{Three Representations, Not Four}
\label{sec:ch10-three-representations-not-four}

\index{representations!the router de-duplicates them|(}%
Four sessions, three statements. The arithmetic is wrong for an N+1, and where
the missing one went is the rest of this section.

\index{representations!what the router actually sends}%
The place to look is the wire, and getting at it means putting something in
front of the Speakers service that writes down each request body as it arrives.
That is an instrument rather than part of the program, so this book does not
print what it saw; the chapter's research note carries the bodies and the
procedure. What they show is that the fetch carries \emph{three} representations
for a list of four sessions, holding the node ids of Ada Fischer, Bruno
Kaminski and Chidi Okafor, in the order the sessions first named them. Two of
those sessions are Ada Fischer's and her identifier is in the list once.

\index{representations!a page of one speaker is one representation}%
Narrow the page to the two sessions that share a speaker and the point is
sharper:

\begin{minted}{graphql}
{ sessions(first: 2) { nodes { title speaker { name } } } }
\end{minted}

\index{representations!one entry, two positions in the answer}%
One representation. Two sessions come back in the response, both carrying a
speaker, and the router asked for her once and put the answer in both places.

\subsection{Which of Them Did That}

\index{DataLoader!both halves look the same from outside}%
There are two candidates and only one experiment separates them. A DataLoader
de-duplicates keys, which is the second thing
chapter~\ref{ch:data-without-n-plus-1} said about it, so a service with one
behind its resolver would show a batch of three whether it had been sent three
or four. The observation above is consistent with either.

\index{representations!the de-duplication survives the loader's removal}%
The version with no loader in it settles it. Three statements, not four, for
four sessions; one statement, not two, for the page of two. There is nothing in
that service keeping track of anything, so whatever collapsed the duplicate
did it before the request was sent.

\index{representations!the resolver de-duplicates nothing}%
And the mirror image, which is the assertion I would want if I were reading
this. Send the loaderless service the same representation three times by hand,
with no router in the path:

\begin{minted}[fontsize=\footnotesize,breakanywhere]{graphql}
query E($r: [_Any!]!) {
  _entities(representations: $r) { __typename ... on Speaker { name } }
}
\end{minted}

\begin{minted}[fontsize=\footnotesize,breakanywhere]{json}
{"r":[{"__typename":"Speaker","id":"U3BlYWtlcjox"},
      {"__typename":"Speaker","id":"U3BlYWtlcjox"},
      {"__typename":"Speaker","id":"U3BlYWtlcjox"}]}
\end{minted}

\index{statement count!three for one row}%
Three statements for one row, and three copies of Ada Fischer in the answer.
The de-duplication is the router's and belongs to the router's entity fetch.
It is not a property of \code{\_entities}, and nothing else calling that route
gets any of it.

\subsection{Where It Happens, and What It Hashes}

\index{graphql-go-tools@\code{graphql-go-tools}!hashes the rendered representation}%
The router's planner and executor are \code{graphql-go-tools}, which the binary
names in its own version output, and the entity fetch it builds is one function.
It hashes each rendered representation, keeps a map from that hash to a position
in the batch it is assembling, and when a hash repeats it appends the original
item to the position already there rather than adding an entry. One answer,
fanned back out to every position that asked for it.

\index{sources!the vendor describing its own router}%
WunderGraph says as much in its own words. Jens Neuse, who co-founded the
company, writes on its blog that the entity route
\enquote{supports a list of representations as input} and that this
\enquote{allows you to batch requests by default}, and, of the router's
algorithm, that \enquote{if there are duplicate representations, we will
deduplicate them before building the variables
object}~\autocite{neuse2023dataloader}. The same post argues against the fix the
rest of this chapter is about, on benchmarks of its own that I have not
reproduced; the two statements above are what this book takes from it.

\index{representations!the hash covers the whole representation}%
The detail worth carrying forward is what goes into the hash: the rendered
representation, not the key inside it. Every representation in this chapter is a
type name and an \code{id} and nothing else, so equal keys mean equal bytes. Two
representations carrying the same key and different \code{@requires} fields
would not be equal bytes. Whether they collapse is unmeasured here, and this
graph has one requiring field, so no query in the book produces the case.

\index{N+1@N+1!N is distinct entities, not rows}%
That fixes the shape of the problem. The N is not the number of rows in the
parent list. It is the number of \emph{distinct} entities in it, which for a
conference schedule where speakers give two talks each is half the rows, and for
a list of orders each belonging to a different customer is all of them. The
router has already done the deduplicating half of what a DataLoader does. The
half it has not done, and cannot do from where it stands, is turn one list into
one query, and that is the half the statements are.
\index{representations!the router de-duplicates them|)}%
